% =====================================================================
%  Chapter 17 — The Proof in Natural Language (new Part IV)
% =====================================================================
\part{The Proof in Natural Language}
\label{part:proof}

\chapter{The Theorem and the Plan of the Proof}
\label{ch:plan}

\begin{seamnote}{Attribution}
Everything in this part is the mathematics of \textbf{Levent Alpöge}. It is
presented here in natural language, so that a reader with a mathematics
background can follow the argument without reading Lean. The corresponding
Lean declarations are cited at the head of each statement, so that every step
can be checked against the formalisation; Chapter~\ref{ch:map} gives the
complete index.
\end{seamnote}

The previous part set out the method step by step, with the source alongside.
This part does the same thing in the other direction: it writes the argument
the way a number theorist would write it, in propositions and proofs, so that
the shape of the mathematics can be judged on its own terms rather than
through the apparatus.

Two conventions apply throughout.

\begin{quote}
\textbf{First:} where a step of the argument is not settled --- and two of
them are not settled, as Chapters~\ref{ch:budget} and~\ref{ch:vanishing} show
--- the gap is marked at the point where it occurs, not deferred to the
end. A proof that hides its gaps is not a proof, and the point of writing
this part is to make the gaps findable.

\textbf{Second:} each proposition carries the Lean declaration that
establishes it, where one exists. That is the honest form of the exercise:
the prose is a reading of the formalisation, and the formalisation is the
checkable object.
\end{quote}

\section{The theorem}\begin{quote}
\textbf{Theorem A (Alp\\\"oge, 2026).} Let
\[
\zeta(s) = \sum_{n \ge 1} n^{-s}\qquad (\mathrm{Re}\,  s > 1)
\]
be analytically continued to the complex plane, meromorphic with a single
simple pole at $s = 1$. Then every zero $\rho$ of $\zeta$ in the critical
strip $0 < \mathrm{Re}\, \rho < 1$ satisfies $\mathrm{Re}\, \rho = \tfrac12$.

Equivalently, and in the form the formalisation uses: for every
Dirichlet character $\chi$ modulo $N$ and every $s$ with $\mathrm{Re}\,  s > \tfrac12$,
\[
L(s,\chi) = \sum_{n \ge 1}\frac{\chi(n)}{n^s}
\]
is non-zero, save for the pole of the principal character at $s = 1$.
\end{quote}

The formal counterpart is
\code{Development/AbelianVarietyTheory/Definitions/Foundations/FinalTheorems.lean},
where the proposition
\code{DirichletLNonvanishingRightOfHalf} is stated and proved, together with
\code{DirichletLZerosOnLine} and
\code{PrimitiveDirichletLZerosOnLine}; and the reduction to the first is
\code{UpperTheory.riemannHypothesis\_of\_everyCharacter} in
\code{CharacterBoxes/RiemannHypothesisOfEveryCharacter.lean}.

\section{Why the statement is the right one}It is worth being explicit about this, because a formalisation can always
prove something and the question is whether it proved \emph{this}.

The equivalence between the two forms of the theorem is standard and was
spelled out in Chapter~\ref{ch:rh}: non-vanishing to the right of the
critical line, together with the functional equation $\Lambda(s) =
\Lambda(1-s)$ and the classification of the trivial zeros, gives RH. The
formalisation performs exactly this reduction, and it is the first thing one
should check in any such development --- because if the reduction is wrong,
nothing else matters.

The formalised statement is not weakened. It is quantified over \emph{all}
characters, not only the primitive ones; the exceptional case
$(\chi = 1, s = 1)$ is a pole and not a zero, and the statement excludes it
for the right reason; and the trivial zeros $0, -2, -4, \dots$ are excluded
by the hypothesis that the zero is not one of them.

\section{The plan}The proof has four blocks and one obstruction.

\begin{enumerate}
\item \textbf{The target.} Reduce RH to a positivity statement about the Weil
expression $W$, using Weil's criterion. This is not a reduction of difficulty;
it is a change of subject, and it is the right one: $W \ge 0$ is a statement
about a quadratic form, and quadratic forms can be certified by matrices.

\item \textbf{The arithmetic factor.} Construct, from the coefficients of a
Hecke character and a quadratic form in $k$ variables, a count whose size is
controlled by a Weil bound. This is where the method departs from its
function-field ancestor and supplies the missing dimension.

\item \textbf{The substitute for multiplicity.} Chain theta lifts to
manufacture the multiplicity that the finite-field method gets from the
independence of the geometric and arithmetic directions, and pay for it with
divisibility.

\item \textbf{The averaging.} Use the Rallis seesaw to remove the character
from the period, so that the sum over characters can be taken, and use
Behrend's construction to control the variance of that sum.

\item \textbf{The obstruction.} Kill the surviving cohomology by twisting.
Here the method requires a result about the $p$-part of the class groups of
cyclotomic fields that is not known.
\end{enumerate}

\section{Notation}Fix a CM field $E$ with maximal real subfield $F$, an odd prime $\ell$, and
a level. Write:

\begin{description}[style=nextline,leftmargin=0pt,font=\sffamily\bfseries\small,
  itemsep=6pt]
\item[$m$, $k$] The weight and the rank of the data. The formalisation's
budget arithmetic requires $k \ge 5$.
\item[$N_s$] The size of the integral structure; $N_s \ge 1$.
\item[$c$] The \emph{available growth}: the multiplicative factor by which the
chain of lifts enlarges the relevant space. The budget requires
$40\,(kN_s) \le c$.
\item[$q$] The \emph{budget}: what the segment of the argument actually
reaches. The budget requires $\tfrac{9}{10}c \le q$.
\item[$W$] The Weil expression, whose non-negativity is equivalent to RH.
\item[$g$] An admissible test function for Weil's criterion.
\end{description}

The last two of these --- $c$ and $q$ --- are the reason the argument has a
ceiling, and they are introduced in Chapter~\ref{ch:budget}. They are named
here because the plan references them.

\section{What a reader should expect}

The proof is not short, and it is not uniformly hard. The honest distribution
is:

\begin{center}
\begin{tabular}{@{}>{\RaggedRight\arraybackslash\hyphenpenalty=10000\exhyphenpenalty=10000}p{0.26\textwidth}>{\RaggedRight\arraybackslash\hyphenpenalty=10000\exhyphenpenalty=10000}p{0.62\textwidth}@{}}
\toprule
\textbf{Block} & \textbf{How hard it is} \\
\midrule
The target & Standard. Weil's criterion is classical and the reduction is a
routine argument with the functional equation. \\
The arithmetic factor & The genuinely new idea. It is a character-sum
construction whose sign comes from the middle cohomology of a quadric, built
so that the same construction is available over the integers. \\
The lift chain & Classical theta-lift machinery, applied in a chain. The
machinery is well understood; the chaining is the claim. \\
The averaging & Classical. Behrend's construction and the seesaw are both
standard. \\
The obstruction & \textbf{Not done.} This requires an open conjecture, or a
weakening of one that is not known. \\
\bottomrule
\end{tabular}
\end{center}

\begin{keytakeaways}
\item The theorem proved is the full Riemann hypothesis, and equivalently the
generalised hypothesis for Dirichlet $L$-functions, with no weakening of the
statement.
\item The proof has four blocks and one obstruction: a positivity target, an
arithmetic factor, a substitute for multiplicity, an averaging, and a final
vanishing step that is not settled.
\item Two parameters, the available growth $c$ and the budget $q$, govern
whether the argument closes, and their relation is where the ceiling comes
from.
\end{keytakeaways}
